\section{A guarded reduction from 240 to 239 states}
\label{sec:guarded-local-239}

Let $M$ be the pinned Rotation240 machine, using its inherited state
labels, start state $95$, and halt state $254$.  Let $N$ be
\path{machines/zfc239_old.tm}.  A configuration is $(q,h,T)$ with
$T:\mathbb Z\to\{0,1\}$; the local statements below allow an arbitrary
exterior tape.  The canonical presentation
\path{machines/zfc239.tm} has working states $0,\ldots,238$, start $0$,
dispatcher $84$, and halt $239$.

Define $\phi(169)=80$ and $\phi(q)=q$ otherwise, and put
$\Phi(q,h,T)=(\phi(q),h,T)$. The exact row changes were given in
Section~\ref{sec:construction}.

The extra freedom comes from reusing a helper in two tape contexts.
The modified row at $91$ remains equivalent to its old behavior in
old state-$91$ contexts, while an entry through the merged role has a
different nearby bit and follows a different return path.  Consequently
the proof concerns old configurations and corresponding new
checkpoints; it does not assert that the old invariant holds at every
intermediate step of $N$.

\subsection{A syntactic invariant and one additional guard}

For an old configuration write $t_i=T(h+i)$.  Consider
\begin{equation}
\label{eq:g239-P}
P(q,h,T):\quad
\begin{cases}
t_{-1}=0,&q=67,\\
(t_{-2},t_{-1})=(0,1),&q=80,\\
(t_{-3},t_{-2},t_{-1})=(0,1,1),&q=91,\\
t_{-1}=0,&q=169.
\end{cases}
\end{equation}
Unlisted states impose no condition.  This is an invariant of $M$
started at $95$ on \emph{every} tape.  Its proof uses the complete
incoming-row inventory:
\begin{align*}
67&:\quad (52,0)\mapsto(0,R,67),\quad (124,0)\mapsto(0,R,67),\\
80&:\quad (67,1)\mapsto(1,R,80),\\
91&:\quad (80,1)\mapsto(1,R,91),\\
169&:\quad (147,0)\mapsto(0,R,169).
\end{align*}
Initialization is vacuous at $95$.  Each incoming write gives the
nearest required bit; the invariant at its predecessor supplies any
older bits.  Induction on the number of old steps proves
\eqref{eq:g239-P}.

One further property is needed:
\begin{equation}
\label{eq:g239-Q}
Q(q,h,T):\quad q=169\ \wedge\ t_0=1\quad\Longrightarrow\quad t_1=1.
\end{equation}
We establish this property for the exact blank initialization and all
valid instruction frames.  We make no all-tape claim for
\eqref{eq:g239-Q}.

\subsection{Four exact positive blocks}

The following paths include both their initial state and final state.
Each new path has exactly the same final tape and head position as
the indicated single old transition, and its final state is mapped
by $\phi$.
\begin{align*}
(80,0),\quad t_{-1}=1:
&\qquad 80\to106\to92\to107\to125\to25,
\tag{5 steps}\label{eq:g239-block80}\\
(169,0),\quad t_{-1}=0:
&\qquad 80\to106\to93\to57,
\tag{3 steps}\label{eq:g239-block169zero}\\
(169,1),\quad t_{-1}=0,\ t_1=1:
&\qquad 80\to91\to167\to151\to173\to56.
\tag{5 steps}\label{eq:g239-block169one}
\end{align*}
For old row $(91,1)$, when
$(t_{-3},t_{-2},t_{-1})=(0,1,1)$, the new path is
\begin{equation}
91\to167\to151\to24\to36\to54\to71\to82\to94\to50.
\tag{9 steps}\label{eq:g239-block91}
\end{equation}
These are literal transition calculations.  The first two paths read
only offsets $-1,0$ and leave their bits unchanged.  The third reads
only $-1,0,1$, changes the original head cell from $1$ to $0$, and
ends one cell right.  The last reads only $-3,-2,-1,0$, changes just
the original head cell from $1$ to $0$, and ends one cell left.
Thus every exterior cell is untouched.  No path contains an internal
halt.

The other $476$ old rows commute with $\Phi$ in one new step.
The program \path{verification/verify_guarded_fold239.py} independently
checks this inventory, invariant closure on $121{,}536$ admitted
Boolean windows, and all $22$ admitted assignments to the common
five-cell footprint of the four blocks.  Its whole-tape comparisons
are exact; cells outside the represented footprint are never read or
written.  The resulting receipt is
\path{evidence/guarded_fold239_certificate.json}.

\subsection{Unbounded certification of the right-hand guard}

The driver \path{verification/certify_guard169.py} instruments the
unchanged affine/list checker \path{verification/symbolic_runtime.py}.
It retains every existing whole-tape endpoint and positive-clock
check.  Register values, additional gap, first spare value, and the
finite remaining spare list stay arbitrary.  The complete pinned
source tree has $6{,}455$ leaves and $8{,}386$ instruction cases after
splitting decrement guards.

At each explicit visit to $169$, the driver normalizes the symbolic
tape and verifies that the head scans rightward into an explicit
run.  If that run contains ones, its length has the form
$c+\sum_i a_i x_i$, with all $a_i,x_i\ge0$, and the driver checks
$c\ge2$.  This proves \eqref{eq:g239-Q} uniformly.  Across the complete
audit there are $32$ such symbolic visits, and the smallest certified
constant is in fact $21$.

There are no hidden visits inside accelerated steps.  Neither row at
$169$ is a self-loop.  State $169$ is neither an entry of a list
lemma nor an interior state in any of the four separately rederived
one-block proofs.  Natural-number induction on unary scans and
structural induction on finite lists therefore preserve the
instrumentation's coverage.  The audit closes all $8{,}386$ cases
using $2{,}322{,}575$ symbolic rules and directly covers all $480$ old
rows.  The proof-construction cap bounds neither tape parameters nor
actual execution lengths.  Receipts and the affine visit records are
\path{evidence/guard169_universal.json} and
\path{evidence/guard169_visits.json}.

The separately checked blank prefix has $409{,}691$ old steps and
six old $(169,1)$ visits, each satisfying the guard.  Its endpoint
is the baseline clean dispatcher frame.  The runtime frame theorem
then supplies the guard throughout continuation.  These are ordinary
mathematical inductions with checked finite premises; this bundle
does not claim a new Lean or Rocq kernel proof.

\subsection{Global simulation and the initialization clock}

For any old execution satisfying $P$ and $Q$, concatenate the
appropriate positive blocks.  If $C_n$ is its $n$th configuration,
the new execution has configuration $\Phi(C_n)$ at time
\begin{align*}
\tau(n)=n
 &+4V_{80,0}(n)+2V_{169,0}(n)\\
 &+4V_{169,1}(n)+8V_{91,1}(n),
\end{align*}
where $V_{q,b}(n)$ counts the first $n$ old transitions using row
$(q,b)$.  In particular $n\le\tau(n)\le9n$.  Induction proves the
checkpoint identity; positivity prevents an infinite old execution
from collapsing to finitely many new steps.  The absence of internal
halts proves halting equivalence in both directions.

The exact blank census is
\[
(V_{80,0},V_{169,0},V_{169,1},V_{91,1})=(18,12,6,3).
\]
It gives an additional $4\cdot18+2\cdot12+4\cdot6+8\cdot3=144$
steps.  Thus the same complete frame occurs after $409{,}835$ new
steps.  The local verifier compares all $409{,}692$ old checkpoints
with their new counterparts, including the entire finite support of
the tape.  Both endpoints have state $96$, head $-5$, and $351$ one
cells, with every other cell zero.

The resulting blank-run equivalence is expressible and provable in
ZFC by the elementary inductions above.  Composing it with the
separately supplied Rotation240 source calibration transfers
$\mathrm{NH}(M)\leftrightarrow\operatorname{Con}(\mathrm{ZFC})$ to
$\mathrm{NH}(N)\leftrightarrow\operatorname{Con}(\mathrm{ZFC})$.
The source calculus, enumerator, and logical calibration are inherited;
the new contribution here is the coordinated state reuse and its
guarded positive-block proof.
